\section*{Bab \ref{ch:g4:animal-reproduction} --- Bagaimana Hewan Berkembang Biak}

\begin{solution}{exo:g4:animal-reproduction:1}
Penggabungan satu \omterm{def:g4:animal-reproduction:sexual}{sel sperma} dengan satu \omterm{def:g4:animal-reproduction:sexual}{sel telur}, yang memulai
seekor \omterm{def:g1:animals-around-us:animal}{hewan} muda yang baru. \omterm{def:g4:animal-reproduction:sexual}{Sel telurnya} datang dari \omterm{def:g4:animal-reproduction:sexual}{betinanya}, \omterm{def:g4:animal-reproduction:sexual}{sel
spermanya} dari \omterm{def:g4:animal-reproduction:sexual}{jantannya}.
\end{solution}

\begin{solution}{exo:g4:animal-reproduction:2}
\omterm{def:g4:animal-reproduction:ovivivi}{Ovipar}: \omterm{def:g4:animal-reproduction:sexual}{betinanya} mengeluarkan \omterm{def:g2:animals-grow-up:born}{telur} dan anaknya menyelesaikan
pertumbuhannya di dalam \omterm{def:g2:animals-grow-up:born}{telur}, di luar tubuhnya --- ayam \omterm{def:g4:animal-reproduction:sexual}{betina}, katak
(juga \omterm{prop:g3:sorting-living-things:groups}{ikan} trout, kadal, siput, \omterm{prop:g3:sorting-living-things:groups}{serangga}). \omterm{def:g4:animal-reproduction:ovivivi}{Vivipar}: anaknya tumbuh di
dalam induknya, diberi makan olehnya, lalu \omterm{def:g2:animals-grow-up:born}{lahir hidup} --- kuda
\omterm{def:g4:animal-reproduction:sexual}{betina}, paus (hampir semua \omterm{prop:g3:sorting-living-things:groups}{mamalia}).
\end{solution}

\begin{solution}{exo:g4:animal-reproduction:3}
Para \omterm{def:g4:animal-reproduction:sexual}{jantan} bernyanyi untuk memanggil para \omterm{def:g4:animal-reproduction:sexual}{betina} ke kolamnya; tiap
pasangan melepaskan selnya ke dalam air, dan di situ sperma bertemu
\omterm{def:g2:animals-grow-up:born}{telur} --- \omterm{def:g4:animal-reproduction:sexual}{pembuahan}; \omterm{def:g2:animals-grow-up:born}{telur} yang terbuahi di dalam agar-agarnya tumbuh
menjadi \omterm{ex:g2:animals-grow-up:frog}{berudu} bulan Juni.
\end{solution}

\begin{solution}{exo:g4:animal-reproduction:4}
\omterm{def:g4:animal-reproduction:sexual}{Pembuahan} \omterm{prop:g3:sorting-living-things:groups}{ikan} troutnya terjadi di dalam air, di luar tubuh \omterm{def:g4:animal-reproduction:sexual}{betinanya}
--- \omterm{def:g2:animals-grow-up:born}{telur} dan sperma dikeluarkan bersama-sama. \omterm{def:g4:animal-reproduction:sexual}{Pembuahan} ayam
\omterm{def:g4:animal-reproduction:sexual}{betinanya} terjadi di dalam tubuhnya, ketika \omterm{prop:g4:animal-reproduction:where}{kawin}. Di darat sel yang
telanjang akan mengering; di dalam air keduanya dapat bertemu dengan
aman di luar tubuh.
\end{solution}

\begin{solution}{exo:g4:animal-reproduction:5}
Anak ayamnya tumbuh di dalam \omterm{def:g2:animals-grow-up:born}{telur} bercangkang di luar tubuh induknya,
hidup dari kuning telur yang dikemas bersamanya. Anak kudanya tumbuh
sebelas bulan di dalam kuda \omterm{def:g4:animal-reproduction:sexual}{betinanya}, diberi makan terus-menerus oleh
tubuhnya, lalu \omterm{def:g2:animals-grow-up:born}{lahir hidup}.
\end{solution}

\begin{solution}{exo:g4:animal-reproduction:6}
Makin banyak anak yang dihasilkan sebuah spesies, makin sedikit
perawatan yang diterima tiap anaknya; makin sedikit anaknya, makin
banyak perawatannya. Ujung-ujungnya: jutaan \omterm{def:g2:animals-grow-up:born}{telur} \omterm{prop:g3:sorting-living-things:groups}{ikan} kod yang tidak
dirawat, dan satu anak gajah yang dijaga selama sepuluh tahun.
\end{solution}

\begin{solution}{exo:g4:animal-reproduction:7}
Karena keberhasilan hanya menuntut bahwa, rata-rata, cukup banyak anak
yang bertahan untuk menggantikan induknya. Dari jutaan \omterm{def:g2:animals-grow-up:born}{telur}, selalu
ada beberapa yang lolos dari bahaya laut --- jumlah yang amat besar
itu \emph{adalah} perlindungannya.
\end{solution}

\begin{solution}{exo:g4:animal-reproduction:8}
Sedikit anak (lima itu sedikit dibandingkan ribuan), tumbuh di dalam
induknya, disusui dan dijaga: kucingnya jelas berada di sisi
sedikit-dan-terawat --- seperti \omterm{prop:g3:sorting-living-things:groups}{mamalia} pada umumnya.
\end{solution}

\begin{solution}{exo:g4:animal-reproduction:9}
Bahwa ia menerima sesuatu dari \emph{tiap} induknya --- anaknya
bukanlah salinan induk \omterm{def:g4:animal-reproduction:sexual}{betinanya} saja. Kedua induknya menyumbang pada
apa yang menjadi anak kudanya.
\end{solution}

\begin{solution}{exo:g4:animal-reproduction:10}
Sebuah kelompok \omterm{def:g2:animals-grow-up:born}{telur} yang kecil --- beberapa lusin atau ratus butir,
bukan jutaan. Perawatan dan jumlah saling menukar: seorang ayah yang
menjaga, mengipasi dan membela dapat membawa kelompok \omterm{def:g2:animals-grow-up:born}{telur} yang kecil
melewati bahaya yang akan menuntut jutaan \omterm{def:g2:animals-grow-up:born}{telur} yang tak terjaga.
\end{solution}

\begin{solution}{exo:g4:animal-reproduction:11}
\omterm{def:g4:animal-reproduction:ovivivi}{Ovipar} --- dan tidak dirawat: kelompok \omterm{def:g2:animals-grow-up:born}{telur} yang terkubur ditinggalkan
kepada kehangatan pasirnya, dan tukiknya berlari sendirian ke laut
(jumlah mereka yang besar itulah siasatnya). Ia harus bertelur di
darat karena \omterm{def:g2:animals-grow-up:born}{telur} \omterm{prop:g4:classifying-animals:classes}{reptil} bercangkang dan bernapas dengan udara ---
kalau dikeluarkan di laut, \omterm{def:g2:animals-grow-up:born}{telurnya} akan tenggelam. Ciri kelasnya,
bukan kebiasaannya, yang menetapkan aturannya.
\end{solution}
